\documentstyle[%


righttag%


,11pt%


,amssymb%


,russian%               remove in Eng.


,izvru%                 Set izven% in Eng.


% ,izvru-ks             for brief communications


]{amsart}





%       Math definitions





\newcommand{\eqdef}{\overset{\mathrm{def}}{=}{}}


\newcommand{\vraisup}{\operatornamewithlimits{vrai\;sup}}


\newcommand{\mes}{\operatorname{mes}}


\newcommand{\const}{\operatorname{const}}


\newcommand{\arctg}{\operatorname{arctg}}


\newcommand{\vraiinf}{\operatornamewithlimits{vrai\;inf}}


\newcommand{\var}{\operatornamewithlimits{var}}


\newcommand{\tts}[1]{}





%\hoffset=-15mm%                  For Eng.


\setcounter{page}{56}


\god{2000}


\nomer{1~(452)} %                        Remove in Eng.


%\nomer{Vol.\,44, No.\,1}%               Set in Eng.


%\str{\pageref{begin}--\pageref{end}}%  Set in Eng.


\udk{517.984.6\,:\,517.929}





\author{\it Л.Ф.\,Рахматуллина}


% NB:   ^^ remove in Eng.


\title{Верхняя оценка спектрального радиуса изотонного оператора}


\thanks{Работа выполнена при финансовой поддержке Российского фонда


фундаментальных исследований (гранты \No\,96-15-96195,


\No\,99-01-01278)


и Конкурсного центра по исследованиям в области фундаментального


естествознания, Санкт-Петербург.}


\date{}


%\pagestyle{myheadings}%         Set in Eng.





\pagestyle{plain} %              Remove in Eng.


\begin{document}


%\thispagestyle{plain}%          Set in Eng.


\maketitle


\label{begin}








Для изотонного интегрального вполне непрерывного оператора


$$


(Ax)(t)\eqdef \int_a^b K(t,s) x(s) ds\quad


(K(t,s)\geq 0, (t,s)\in[a,b]\times[a,b])


$$


в пространстве $C[a,b]$ непрерывных на $[a,b]$


функций справедливо утверждение: спектральный радиус


$\rho(A)$ оператора $A$


меньше единицы тогда и только тогда, когда выполнено условие A:


существует такая непрерывная функция $v$, что


$$


v(t)\geq 0,\quad r(t)\eqdef v(t) - (Av)(t) \geq 0,\quad t\in[a,b],


$$


причем множество нулей невязки $r$


не более чем счетно. Это утверждение играет существенную роль в теории


линейных дифференциальных уравнений. В теории 


функционально-дифференциальных уравнений


возникает потребность в установлении оценки $\rho(A)<1$


для изотонного оператора


$A:C[a,b]\to C[a,b]$,


который не является интегральным


([1], с.\,133).


Приведенное выше утверждение является частным случаем теоремы Г.Г.\,Исламова


[2], [3], согласно которой для изотонного слабо компактного оператора


$A:C[a,b]\to C[a,b]$ неравенство $\rho(A)<1$


справедливо тогда и только тогда, когда выполняется условие A,


причем $r(t) >0$ в специальных точках отрезка


$[a,b],$


называемых ``существенными точками''. Ослабить условие о нулях невязки $r$


и отказаться от требования слабой компактности удается за счет специальных


свойств


оператора [4]. Ниже предлагается дальнейшее развитие идей работы [4].





Пусть $T\subset R^1$ --- измеримое по Лебегу множество,


$\mes T\leq + \infty$. Через $C$


обозначим банахово пространство ограниченных непрерывных функций


$x:T\to R^1$ с нормой


$\|x\|_C=\sup\limits_{t\in T} |x(t)|.$ Пусть


$\gamma:T\to R^1$ --- непрерывная функция,


$\gamma(t)>0$, $t\in T$;


$C^\gamma$\ ---


банахово пространство таких непрерывных функций $x: T\to R^1$,


что


$$


\frac{x}{\gamma}\in C;\quad \|x\|_{C^{\gamma}} = \sup_{t\in


T}\frac{|x(t)|}{\gamma(t)}.


$$


Через $L_\infty$


обозначим банахово пространство измеримых ограниченных в существенном


функций $x:T\to R^1$,


$\|x\|_{L_{\infty}} = \vraisup\limits_{t\in T}|x(t)|$.


Если $\gamma:T\to R^1$ --- измеримая функция, $\gamma(t) >0$


почти всюду на $T$, то $L_\infty^\gamma$ ---


банахово пространство таких измеримых функций $x: T\to R^1$, что


$\frac{x}{\gamma}\in L_\infty$,


$\|x\|_{L_{\infty}^{\gamma}} = \vraisup\limits_{t\in T}


\frac{|x(t)|}{\gamma(t)}$.





В пространствах $C^\gamma$ и $L_\infty^\gamma$


естественным образом вводится полуупорядоченность. А именно, для


$x\in C^\gamma$ полагают $x>0$, если $x(t)\geq 0$ для любого


$t\in T$ и $x(t)\not\equiv 0$ на $T$; $x=0$, если $x(t)\equiv 0$


на $T$; $x_1 <x_2$, если $x_2-x_1>0$. Аналогично для


$L_\infty^\gamma$: $x>0$, если $x(t)\geq 0$ почти всюду на $T$ и


$x(t) >0$ на множестве положительной меры; $x=0$, если


$x(t) = 0$ почти всюду на $T$; $x_1<x_2$, если $x_2-x_1>0$.





Ниже символом $B^\gamma$ будем обозначать одно из пространств


$C^\gamma$ и $L_\infty^\gamma$.


При этом в формулировках и доказательствах утверждений в случае пространства


$C^\gamma$ символы ``$\vraisup$'' и ``$\vraiinf$''


заменяются на эквивалентные им символы ``$\sup$'' и ``$\inf$''.


Естественно, что в некоторых случаях доказательства упрощаются.


\medskip





{\bf Определение.} Линейный оператор $A:B^\gamma\to B^\gamma$


называется изотонным, если для любого $x\in B^{\gamma}$


из $x>0$ следует $Ax\geq 0$. Отсюда $Ax_1\leq Ax_2$, если $x_1<x_2$.


\medskip





Ниже будет использована норма в пространстве $B^\gamma$,


отличная от нормы


$\|\cdot\|_{B^{\gamma}}$,


но эквивалентная ей. А именно, пусть $v\in B^\gamma$,


$\vraiinf\limits_{t\in T}\frac{v(t)}{\gamma(t)} >0$. Положим


$\|x\|_v = \vraisup\limits_{t\in


T}\frac{|x(t)|}{v(t)}$ и обозначим через


$\|A\|_v$ норму оператора $A:B^\gamma\to B^\gamma$, соответствующую норме


$\|\cdot\|_v$.


\medskip





{\bf Утверждение.} {\it


Пусть $A:B^\gamma\to B^\gamma$\ ---


линейный ограниченный изотонный оператор. Тогда


$\|A\|_v = \|Av\|_v$.


}


\medskip





{\bf Доказательство.} Имеем


$\|Av\|_v\leq  \|A\|_v \|v\|_v = \|A\|_v$. Пусть


$\|x\|_v=1$. Тогда для любого $\varepsilon >0$


существует такое множество $U\subset T$, $\mes U=0$, что


$\frac{|x(t)|}{v(t)}<1+\varepsilon$ для $t\in T\setminus U$.


Следовательно, для $\|x\|_v=1$


$|(Ax)(t)|\leq(1+\varepsilon)(Av)(t)$ почти всюду на $T$.


Отсюда


\begin{gather*}


\vraisup_{t\in T}\frac{|(Ax)(t)|}{v(t)} \leq (1+\varepsilon)


\vraisup_{t\in T}\frac{(Av)(t)}{v(t)},\\


%


\|A\|_v =


\sup_{\|x\|_v=1} \|Ax\|_v=


\sup_{\|x\|_v=1}


\bigg\{\vraisup_{t\in T}\frac{|(Ax)(t)|}{v(t)}\bigg\}


\leq (1+\varepsilon)\|Av\|_v.


\end{gather*}


Ввиду произвольности $\varepsilon$\;


$\|A\|_v\leq \|Av\|_v.\quad\square$


\medskip





{\bf Лемма.} {\it


Пусть $A:B^\gamma\to B^\gamma$ ---


линейный ограниченный изотонный оператор.


Неравенство $\rho(A)<1$ справедливо


тогда и только тогда, когда существует такой элемент


$v\in B^\gamma$, что


$$


v>0,\quad \beta = \vraiinf_{t\in T}\frac{v(t)-(Av)(t)}{\gamma(t)}>0.


$$


}





Отметим, что при $T=[a,b]$, $\gamma(t)\equiv 1$


это утверждение хорошо известно.





{\bf Доказательство. Необходимость.} В качестве $v$


можно взять решение уравнения $x-Ax=\gamma$.


Тогда


$v=\gamma + A\gamma + A^2\gamma + \dotsc$,


$v>0$, $\frac{1}{\gamma(t)}[v(t)-(Av)(t)] = 1$


почти всюду на $T$.


\smallskip





{\bf Достаточность.} Воспользуемся нормой $\|\cdot\|_v$.


Имеем


$$


\|A\|_v = \|Av\|_v = \vraisup_{t\in T}\frac{(Av)(t)}{v(t)}.


$$


Зафиксируем произвольно $\beta_1$, $0<\beta_1<\beta$,


и выберем такое множество $U\subset T$, $\mes U=0$, что


$\frac{1}{\gamma(t)}(v(t)-(Av)(t)) >\beta_1$ для


$t\in T\setminus U$ и


$\sup\limits_{t\in T\setminus U}\frac{v(t)}{\gamma(t)} <\infty$.


Тогда


$\frac{(Av)(t)}{v(t)} < 1- \frac{\beta_1}{\frac{v(t)}{\gamma(t)}}$


для $t\in T\setminus U$. Отсюда


$$


\|Av\|_v \leq \sup_{t\in T\setminus U} \frac{(Av)(t)}{v(t)} \leq


1 - \frac{\beta_1}{\sup\limits_{t\in T\setminus U}\frac{v(t)}{\gamma(t)}}<1.


$$


Следовательно, $\rho(A)<1.\quad\square$





Требования относительно $v$ и $r=v-Av$


можно ослабить за счет дополнительных ограничений на оператор $A$.


Одним из таких ограничений (естественных при рассмотрении ряда краевых


задач) является свойство $\cal M$.





Будем говорить, что линейный оператор $A:B^\gamma\to B^\gamma$


обладает {\it свойством} $\cal M$, если


$$


\vraiinf_{t\in T} \frac{(Ax)(t)}{\gamma(t)} >0


$$


для любого $x\in B^\gamma$, $x>0$.





\begin{thm}


Пусть линейный ограниченный оператор $A:B^\gamma\to B^\gamma$


обладает свойством $\cal M$. Если существует такой элемент


$v\in B^\gamma$, что


$v>0$, $r\eqdef v-Av >0$,


то $\rho(A)<1$.


\end{thm}





{\bf Доказательство.}


Если


$\vraiinf\limits_{t\in T} \frac{r(t)}{\gamma(t)} >0$,


то $\rho(A)<1$ в силу леммы.





Пусть $\vraiinf\limits_{t\in T} \frac{r(t)}{\gamma(t)} =0$.


Применяя оператор $A$ к обеим частям равенства $v-Av=r$,


получим $Av-A^2v=Ar$. Отсюда и из неравенства $v-Av>0$ имеем


$$


r_1(t)\eqdef v(t) - (A^2v)(t) \geq (Ar)(t)


$$


почти всюду на $T$. Поэтому


$$


\vraiinf_{t\in T}\frac{r_1(t)}{\gamma(t)}\geq \vraiinf_{t\in T}


\frac{(Ar)(t)}{\gamma(t)} >0.


$$


В силу леммы $\rho(A^2)<1$, отсюда $\rho(A)<1.\quad\square$





Будем говорить, что линейный оператор $A:B^\gamma\to B^\gamma$


обладает {\it свойством} $\cal N$,


если существует такое измеримое множество $\Delta\subset T$


и такой элемент $\varphi\in B^\gamma$, что


$$


\varphi >0,\quad


\vraiinf_{t\in\Delta}\frac{\varphi(t)-(A\varphi)(t)}{\gamma(t)} >0.


$$





\begin{thm}


Пусть линейный ограниченный изотонный оператор


$A:B^\gamma\to B^\gamma$ обладает свойством $\cal N$.


Если существует такой элемент $v\in B^{\gamma}$,


что


$v{>}0$, $r\eqdef v-Av{ >}0$, $\vraiinf\limits_{t\in


T\setminus\Delta}\frac{r(t)}{\gamma(t)}{ >}0$,


то $\rho(A)<1$.


\end{thm}





{\bf Доказательство.}


Если $\mes(T\setminus\Delta)=0$, то выполнено условие леммы.





Пусть $\mes(T\setminus\Delta)>0$. Для любого $\varepsilon>0$\;


$v_\varepsilon\eqdef v+\varepsilon \varphi>0$. Положим


$r_{\varepsilon}\eqdef v_\varepsilon - A v_\varepsilon


= r + \varepsilon \psi$, где


$\psi=\varphi - A\varphi$. В силу свойства $\cal N$


\begin{equation}


\vraiinf_{t\in\Delta} \frac{r_{\varepsilon}(t)}{\gamma(t)} >0.


\end{equation}





Обозначим


$\omega^+=\{t\in T\setminus\Delta:\psi(t)\geq 0\}$,


$\omega^-=\{t\in T\setminus\Delta:\psi(t)< 0\}$


($\omega^+$ и $\omega^-$ определяются с точностью до множества нулевой меры).





Пусть $\mes\omega^+>0$. Так как при $t\in\omega^+$\;


$r_{\varepsilon}(t) \geq r(t)$, то


\begin{equation}


\vraiinf_{t\in\omega^+} \frac{r_{\varepsilon}(t)}{\gamma(t)} >0


\end{equation}


для любого $\varepsilon>0$. Если $\mes\omega^- =0$,


то отсюда с учетом (1)


$\vraiinf\limits_{t\in T} \frac{r_{\varepsilon}(t)}{\gamma(t)} >0$


для любого $\varepsilon>0$. Если $\mes\omega^- >0$ и


$m_1=\vraiinf\limits_{t\in\omega^-} \frac{\psi(t)}{\gamma(t)}$,


$m_2 =\vraiinf\limits_{t\in\omega^-} \frac{r(t)}{\gamma(t)}$, то для


$\varepsilon\in (0,-m_2/m_1)$\;


$\vraiinf\limits_{t\in\omega^-} \frac{r_{\varepsilon}(t)}{\gamma(t)}


\geq m_2 + \varepsilon m_1 >0$


и с учетом (1) и (2)\;


$\vraiinf\limits_{t\in T} \frac{r_{\varepsilon}(t)}{\gamma(t)} >0$.


Если $\mes\omega^+ =0$, то


$\vraiinf\limits_{t\in T\setminus\Delta}


\frac{r_{\varepsilon}(t)}{\gamma(t)} =


\vraiinf\limits_{t\in\omega^-} \frac{r_{\varepsilon}(t)}{\gamma(t)} >0$


для


$\varepsilon\in(0,-m_2/m_1)$. Отсюда и из (1)\;


$\vraiinf\limits_{t\in T} \frac{r_{\varepsilon}(t)}{\gamma(t)} >0$.








Таким образом, при достаточно малых


$\varepsilon>0$\; $v_\varepsilon$


удовлетворяет условиям леммы. Следовательно, $\rho(A)<1.\quad\square$





Будем говорить, что линейный оператор $A:B^\gamma\to B^\gamma$


обладает {\it свойством} $\cal M\cal N$,


если он обладает свойством $\cal N$ и


$\vraiinf\limits_{t\in T\setminus\Delta}


\frac{(Ax)(t)}{\gamma(t)} >0$ для любого $x\in B^\gamma$, $x>0$.





\begin{thm}


Пусть линейный ограниченный изотонный оператор $A:B^\gamma\to B^\gamma$


обладает свойством $\cal M\cal N$, причем


$\varphi-A\varphi >0$. Если существует такой элемент


$v\in B^\gamma$, что


$v>0$,


$\vraiinf\limits_{t\in T\setminus\Delta} \frac{v(t)}{\gamma(t)} >0;$


$r\eqdef v-Av >0$,


то $\rho(A)<1$.


\end{thm}





{\bf Доказательство.} Если


$\vraiinf\limits_{t\in T} \frac{r(t)}{\gamma(t)} >0$,


то $\rho(A)<1$


в силу теоремы 2. В противном случае, повторяя схему доказательства теоремы 1


с заменой $T$ на $T\setminus\Delta$, в силу теоремы 2 получим


$\rho(A^2)<1$. Следовательно, $\rho(A)<1.\quad\square$





\begin{note}


В силу леммы условия теорем 1, 2 и 3 относительно $v$ и $r$


необходимы для выполнения неравенства $\rho(A)<1$.


\end{note}





Если $T$


--- замкнутое ограниченное множество, то в случае пространства


$C^\gamma$ можно положить $\gamma(t)\equiv 1$.


\medskip





{\bf Следствие} из теоремы 2 [4].


Пусть $T=[a,b]$


и для линейного ограниченного изотонного оператора


$A:C[a,b]\to C[a,b]$ выполнено условие: существуют такие


$t_1,\dotsc,t_k\in[a,b]$, что $(Ax)(t_i)=0$,


$i=1,\dotsc,k$, для любого $x\in C[a,b]$.


Справедливо соотношение $\rho(A)<1$


тогда и только тогда, когда существует такой элемент $v\in C[a,b]$,


что $v(t)>0$ и $r(t)>0$ при


$t\in[a,b]\setminus\{t_1,\dotsc,t_k\}$.


\medskip





В условиях следствия оператор $A$ обладает свойством $\cal N$.


Чтобы убедиться в этом, достаточно взять в качестве $\Delta$


объединение таких окрестностей точек $t_1,\dotsc,t_k$,


что для $\varphi(t)\equiv 1$\; $(A\varphi)(t)\leq q<1$,


$t\in\Delta$.





\begin{note}


Если удовлетворяющий условиям следствия оператор $A$


слабо компактен, то утверждение может быть получено


на основе теоремы Г.Г.\,Исламова [3].


\end{note}





\begin{exam}


Рассмотрим краевую задачу


\begin{equation}


\begin{split}


({\cal L}x)(t)


\eqdef t(1-t)\Ddot{x}(t) - \int_0^1 &x(s)\,d_sR(t,s) = f(t),\quad


t\in[0,1],\\


%


x(0)&=x(1)=0


\end{split}


\end{equation}


в предположениях: $R(t,\cdot)$ при почти каждом $t\in[0,1]$


не возрастает на $[0,1]$; $R(\cdot,s)$ при каждом $s\in[0,1]$


и $f(\cdot)$ суммируемы на $[0,1]$.


Следуя [5], [6], решением уравнения $({\cal L}x)(t)=f(t)$


называем удовлетворяющую ему почти всюду на $[0,1]$ функцию


$x:[0,1]\to R^1$, которая обладает свойствами


\begin{itemize}


\item[1)] $x(t)$ непрерывна на $[0,1]$;


\item[2)] производная $\Dot{x}(t)$


абсолютно непрерывна на каждом замкнутом отрезке, содержащемся в интервале


$(0,1)$;


\item[3)] произведение $t(1-t)\Ddot{x}(t)$ суммируемо на $[0,1]$.


\end{itemize}


Линейное пространство функций, удовлетворяющих условиям


1)--3), обозначим через $D$.





Краевая задача


$$


t(1-t)\Ddot{x}(t)=z(t),\quad x(0)=x(1)=0


$$


для любой суммируемой функции $z$ имеет единственное решение


$x\in D$ [6]


$$


x(t)=\int_0^1W(t,s) z(s)\,ds,


$$


где


$$


W(t,s)=


\begin{cases}


-\dfrac{1-t}{1-s},& 0\leq s\leq t\leq 1;\\[3mm]


-\dfrac{t}{s},& 0\leq t<s\leq 1.


\end{cases}


$$





Определим оператор $A:C[0,1]\to C[0,1]$ равенством


$$


(Ax)(t) = \int_0^1 W(t,s) \int_0^1 x(\tau)\,d_{\tau}R(s,\tau)\,ds


$$


и обозначим $g(t)=\int\limits_0^1W(t,s)f(s)\,ds$.


Задача (3) эквивалентна уравнению


$$


x=Ax+g


$$


в пространстве $C[0,1]$. Это следует из того, что $g\in D$ и $Ax\in D$


для любого $x\in C[0,1]$. Таким образом, неравенство $\rho(A)<1$


гарантирует однозначную разрешимость задачи (3) для любой суммируемой


$f$.





Для оценки спектрального радиуса оператора $A$


воспользуемся сначала оценкой его нормы без предположения о монотонности


$R(\cdot,\cdot)$


по второму аргументу при условии, что


$R(s)=\var\limits_{\tau\in[0,1]}R(s,\tau)$ суммируема на


$[0,1]$. Так как


$$


\|A\|\leq \max_{t\in[0,1]}\bigg\{- \int_0^1 W(t,s)R(s)\,ds\bigg\},


$$


то $\rho(A)\leq\|A\|<1$, если


\begin{equation}


(1-t)\int_0^t\frac{R(s)}{1-s}\,ds + t\int_t^1\frac{R(s)}{s}\,ds <1,\quad


t\in[0,1].


\end{equation}


Это условие однозначной разрешимости задачи (3) получено в [6]


в предположении, что $R(t,\cdot)$ не возрастает на $[0,1]$.





В предположении о невозрастании $R(t,\cdot)$ оператор


$A:C[0,1]\to C[0,1]$ изотонный, причем $(Ax)(0)=(Ax)(1)=0$


для любого $x\in C[0,1]$.


Воспользуемся теперь следствием из теоремы 2. Положим


$$


v(t)= - (1-t)\ln\,(1-t) - t\ln t = - \int_0^1 W(t,s)\,ds, \quad


t\in[0,1]


$$


(считаем $t\ln t|_{t=0}=0$). Тогда


$$


r(t) = - \int_0^1 W(t,s)\bigg\{1+\int_0^1 v(\tau)


d_{\tau}R(s,\tau)\bigg\}\,ds.


$$


Таким образом, если почти всюду на $[0,1]$


\begin{equation}


-\int_0^1[(1-\tau)\ln\,(1-\tau) + \tau \ln\tau] d_{\tau}R(t,\tau)\leq 1,


\end{equation}


причем на множестве положительной меры неравенство строгое, то $r(t)>0$,


$t\in(0,1)$, и, следовательно, $\rho(A)<1$.





Решение задачи (3) имеет представление [6]


$$


x(t)=\int_0^1G(t,s)f(s)\,ds,


$$


где $G(t,s)$ --- функция Грина этой задачи. Из равенства


$$


\int_0^1 G(t,s)f(s)\,ds = g(t) + (Ag)(t) + (A^2g)(t) + \dotsb


$$


следует, что решение задачи (3) не принимает положительных значений, если


$f(t)\geq 0$.


Таким образом, каждое из неравенств (4) и (5) гарантирует неравенство


$G(t,s)\leq 0$ в квадрате $[0,1]\times[0,1]$.





В случае уравнения с сосредоточенным отклонением аргумента


\begin{gather*}


t(1-t)\Ddot{x}(t) - p(t)x[h(t)]= f(t),\quad t\in[0,1],\\


%


x(\xi)=0,\ \;\text{если}\ \; \xi\notin[0,1],


\end{gather*}


в предположениях, что $h$ измерима, $p$, $f$ суммируемы на $[0,1]$,


$p(t)\leq 0$, неравенства (4) и (5) принимают вид соответственно


\begin{equation}


-(1-t)\int_0^t\frac{p(s)\sigma(s)}{1-s}\,ds -


t\int_t^1\frac{p(s)\sigma(s)}{s}\,ds< 1,


\end{equation}


где $\sigma(t)$ --- характеристическая функция множества


$\omega=\{t\in[0,1]:h(t)\in[0,1]\}$,


\begin{equation}


\sigma(t)p(t)\{(1-h(t))\ln\,(1-h(t))+h(t)\ln h(t)\} \leq 1


\end{equation}


почти всюду на $[0,1]$, причем, если


$\mes\{[0,1]\setminus\omega\}=0$,


то неравенство строгое на множестве положительной меры.


В последнем неравенстве в случае $h(t)\notin[0,1]$


функцию в фигурных скобках можно считать определенной произвольно.





В некоторых случаях условие (7) оказывается менее ограничительным, чем (6).


Так, для $p(t)=\const$, $h(t)=\theta t$,


$0<\theta<1/2$,


условие (7) дает неравенство


$$


-p<\frac{1}{\max\limits_{t\in[0,1]}v(\theta t)}.


$$


Из (6) получаем


$$


-p<\frac{1}{\max\limits_{t\in[0,1]}v(t)}= 1/\ln 2 <


\frac{1}{\max\limits_{t\in[0,1]}v(\theta t)}.


$$





Условия (5) и (7) однозначной разрешимости задачи (3) уточняют условия,


полученные для уравнения второго порядка в [6].


\end{exam}





\begin{exam}


Рассмотрим краевую задачу


\begin{gather}


({\cal L}x)(t)\eqdef


\Ddot{x}(t) - p(t)\{\Dot{x}[h(t)]+x[h(t)]\} =


f(t),\quad t\in[a,b],\notag \\


%


x(\xi)=\Dot{x}(\xi)=0,\ \;\text{если}\ \;\xi\notin[a,b]\\


%


x(a)=\Dot{x}(b)=0 \notag


\end{gather}


в предположениях: $f$ суммируема, $h$ измерима на $[a,b]$,


$p\in L_\infty$, $p\sigma>0$, где $\sigma(t)$


--- характеристическая функция множества


$\omega=\{t\in[a,b]:h(t)\in[a,b]\}$.


Решением уравнения $({\cal L}x)(t)=f(t)$ называется функция


$x:[a,b]\to R^1$, имеющая абсолютно непрерывную производную


$\Dot{x}(t)$ и удовлетворяющая этому уравнению почти всюду на $[a,b]$.





Подстановка $x(t) = \int\limits_a^b W(t,s) z(s)\,ds,$


где


$$


W(t,s) =


\begin{cases}


(t-s) - (t-a),& a\leq s\leq t\leq b;\\


-(t-a),&a\leq t < s \leq b,


\end{cases}


$$


--- функция Грина задачи $\Ddot{x}(t)=z(t)$, $x(a)=\Dot{x}(b)=0$,


устанавливает взаимно однозначное соответствие между множеством решений


$x$ задачи (8) и множеством решений $z$ уравнения


\begin{equation}


z(t)= p(t) \sigma(t)\int_a^b K(t,s) z(s)\,ds + f(t).


\end{equation}


Здесь


$$


K(t,s)=


\begin{cases}


s-a,& a\leq s\leq h(t)\leq b,\\


1+h(t) -a,& a\leq h(t)<s\leq b,\\


0&\text{в остальных точках.}


\end{cases}


$$


Отметим, что наличие первой производной в уравнении не позволяет


эффективно свести задачу (8) к эквивалентному ей уравнению в пространстве


$C[a,b]$.





Обозначим


\begin{equation}


(Az)(t)=p(t) \sigma(t) \int_a^b K(t,s)z(s)\,ds.


\end{equation}


Рассмотрим сначала оператор $A$,


определяемый равенством (10), как оператор, действующий в пространстве


$L_\infty$. Оценивая норму этого оператора (без предположения о знаке


$p$), получим $\|A\|<1$, если


$$


\vraisup_{t\in[a,b]} |p(t)|\sigma(t) \int_a^b K(t,s)\,ds <1.


$$


Для $t\in\omega$


$$


\int_a^b K(t,s)\,ds = \frac{1}{2} (h(t)-a)^2 + (1+h(t)-a)(b-h(t))


\eqdef \psi(t).


$$


Таким образом, неравенство


\begin{equation}


\vraisup_{t\in[a,b]}\{\sigma(t) |p(t)|


[\tfrac{1}{2} (h(t)-a)^2 + (b-h(t))(1+h(t)-a)]\} <1


\end{equation}


гарантирует оценку $\rho(A)\leq \|A\|<1$.





Оператор $A$ при $p(t)\geq 0$ изотонен. Воспользуемся теперь для


оценки его спектрального радиуса теоремой 2. Положим


$\varphi(t)\equiv 1$, тогда


$$


(A\varphi)(t) = p(t)\sigma(t)\int_a^bK(t,s)\,ds = p(t) \sigma(t)\psi(t).


$$


Если $t\in\omega$, то $\psi(t)>0$. Зафиксируем число $q$,


$0<q<1$, и положим


$\Delta=\{t\in[a,b]:0\leq p(t)\sigma(t)\psi(t)\leq q\}$.





Рассмотрим возможные случаи.





1) Пусть $0< \mes\Delta<b-a$. Тогда оператор $A$


обладает свойством $\cal N$


$$


\vraisup_{t\in\Delta} \{p(t)\sigma(t) \psi(t)\}\leq q.


$$


Следовательно,


$$


\vraiinf_{t\in\Delta}\{1-p(t)\sigma(t)\psi(t)\}\geq 1-q>0.


$$





Положим $v(t)=p(t)\sigma(t)$. Тогда


$$


r(t)= p(t)\sigma(t)\bigg\{1-\int_a^bK(t,s)p(s)\sigma(s)\,ds\bigg\}.


$$


В силу выбора множества $\Delta$


$$


\vraiinf_{t\in[a,b]\setminus\Delta} \{p(t)\sigma(t)\}>0.


$$


Так как $[a,b]\setminus\Delta\subset\omega$, то


$$


\vraiinf_{t\in[a,b]\setminus\Delta} r(t)


\geq \vraiinf_{t\in[a,b]\setminus\Delta} \{p(t)\sigma(t)\}


\vraiinf_{t\in\omega} \bigg\{1-


\int_a^bK(t,s)p(s)\sigma(s)\,ds\bigg\}.


$$





Таким образом, если выполнено неравенство


\begin{equation}


\vraisup_{t\in[a,b]} \sigma(t)


\bigg\{\int_a^{h(t)} (s-a)p(s)\sigma(s)\,ds + (1+h(t)-a)


\int_{h(t)}^b p(s)\sigma(s)\,ds\bigg\} <1,


\end{equation}


то $\rho(A)<1$ в силу теоремы 2 (здесь $p(s)$ для $s\notin[a,b]$


следует считать доопределенной произвольно).





2) Если $\mes\Delta=0$, то $p(t)\sigma(t)\psi(t)\geq q$


почти всюду на $[a,b]$. Следовательно, $\sigma(t)=1$ почти всюду на


$[a,b]$ и $\vraiinf\limits_{t\in[a,b]} p(t)>0$.


Положим $v(t)=p(t)$. Тогда


$$


r(t) = p(t)\bigg\{ 1-\int_a^bK(t,s)p(s)\,ds\bigg\}


$$


и в силу леммы $\rho(A)<1$, если


$$


\vraiinf_{t\in[a,b]} \bigg\{1-\int_a^b K(t,s)p(s)\,ds\bigg\}>0,


$$


т.\,е. если выполнено неравенство (12).





3) Если $\mes\Delta=b-a$, то


$\vraisup\limits_{t\in[a,b]} \{p(t)\sigma(t)\psi(t)\}\leq q$


и $\rho(A)<1$ в силу (11).





Так как $L_\infty\subset L_1$, $L_1$ --- пространство суммируемых на


$[a,b]$ функций, и для $z\in L_\infty$


$$


\|z\|_{L_1}= \int_a^b|z(t)|\,dt\leq (b-a)\|z\|_{L_{\infty}},


$$


то спектральный радиус оператора $A$


как оператора, действующего в пространстве суммируемых функций, также


меньше единицы ([7], с.\,75).





Итак, если выполнено неравенство (12), то $\rho(A)<1$


и, следовательно, уравнение (9) и задача (8) однозначно разрешимы для любого


$f\in L_1$. Решение $z$ уравнения (9) имеет представление


$$


z= f+Af+A^2f+\dotsb


$$


Для решения $x$ задачи (8) имеем


$$


x(t)=\int_a^bG(t,s)f(s)\,ds = \int_a^b W(t,s)z(s)\,ds.


$$


Таким образом, если $f>0$, то $x<0$ и, следовательно, для функции Грина


$G(t,s)$ задачи (8) выполняется неравенство $G(t,s)\leq 0$


в квадрате $[a,b]\times[a,b]$.


\end{exam}





\begin{exam}


Задача Коши


$$


\Dot{x}(t)- \frac{2t}{t^2+1} x(t) = z(t),\ \; t\in T=[0,\infty),\ \; x(0)=0


$$


для каждого $x\in L_\infty$ имеет единственное решение


$$


x(t) = (t^2+1)\int_0^t\frac{z(s)}{s^2+1}\,ds,


$$


принадлежащее пространству $C^\gamma$, где


$\gamma(t)=t^2+1$.


Сформулируем условия, при которых задача Коши для уравнения с


отклоняющимся аргументом


\begin{equation}


\Dot{x}(t)- \frac{2t}{t^2+1} x(t) - p(t) x_h(t) = f(t),\ \;


t\in T,\ \; x(0)=0


\end{equation}


также имеет единственное принадлежащее пространству $C^\gamma$


решение для любого $f\in L_\infty$. Здесь $h:[0,\infty)\to R^1$


--- измеримая ограниченная в существенном на каждом конечном


промежутке функция; $p(t)\sigma(t)\geq 0$, $t\in[0,\infty)$,


$\sigma(t)$ --- характеристическая функция множества


$\{t\in[0,\infty):h(t)\geq 0\}$;


$$


x_h(t) =


\begin{cases}


x[h(t)],&\text{если }\; h(t)\geq 0;\\


0,&\text{если }\; h(t)<0.


\end{cases}


$$





Задача (13) эквивалентна уравнению $x=Ax+g$, где


$$


(Ax)(t) = (t^2+1)\int_0^t \frac{p(s)}{s^2+1} x_h(s)\,ds,\quad


g(t)=(t^2+1)\int_0^t\frac{f(s)}{s^2+1}\,ds.


$$


Оператор $A$ действует в пространстве $C^\gamma$


и непрерывен, если


\begin{equation}


\vraisup_{t\in[0,\infty)}\frac{h^2(t)+1}{t^2+1} \sigma(t) = \nu <\infty


\end{equation}


и


\begin{equation}


\int_0^{\infty} p(s)\sigma(s)\,ds <\infty.


\end{equation}


При этом $\|A\|\leq \nu\int\limits_0^{\infty} p(s)\sigma(s)\,ds$.


Отметим, что условие (14) выполняется, например, для функции


$h(t)= k t^{\alpha} +\tau$, где $k,\tau,\alpha = \const$,


$k>0$, $0<\alpha\leq 1$.





Оператор $A:C^{\gamma}\to C^{\gamma}$ изотонен. Он обладает свойством


$\cal N$. Действительно, для $\varphi(t)=t^2+1$ и


$\Delta=[0,\delta]$, где $\delta>0$ выбрано так, что


$\int\limits_0^{\delta} p(s)\sigma(s)\,ds <\frac{1}{\nu}$,


имеем


$$


\min_{t\in[0,\delta]}\frac{\varphi(t)-(A\varphi)(t)}{t^2+1}


\geq 1-\nu\int_0^{\delta} p(s)\sigma(s)\,ds >0.


$$





Положим


$$


v(t) = (t^2+1)\arctg t=(t^2+1)\int_0^t\frac{ds}{s^2+1}.


$$


Имеем


$$


\inf_{t\in(\delta,\infty)}\frac{v(t)}{t^2+1} = \arctg\delta >0,\quad


\frac{r(t)}{t^2+1}=


\int_0^t\frac{1}{s^2+1}\{1-p(s)\sigma(s) v[h(s)]\}\,ds.


$$





Пусть почти всюду на $[0,\infty)$ выполнено неравенство


\begin{equation}


p(s)\sigma(s) [h^2(s)+1]\arctg h(s) \leq 1,


\end{equation}


причем это неравенство строгое на множестве положительной меры, содержащемся в


$[0,\delta]$. Тогда


$\frac{r(t)}{t^2+1}\geq 0$, $t\in[0,\infty)$,


$$


\inf_{t\in(\delta,\infty)}\frac{r(t)}{t^2+1} = \int_0^\delta


\frac{1}{s^2+1} \{ 1- p(s)\sigma(s)


v[h(s)]\}\,ds >0.


$$





Итак, если выполнены (14)--(16), то $\rho(A)<1$ в силу теоремы 2


и, следовательно, уравнение $x=Ax+g$


и задача (13) имеют единственное решение


$x\in C^\gamma$ при любом $f\in L_\infty$.


\end{exam}





\begin{note}


Оценка (16) иногда оказывается менее ограничительной, чем условие,


получаемое на основе неравенства $\|A\|<1$. Так, для $h(t)=t$ и


\begin{gather*}


p(t)=


\begin{cases}


0,&\text{если }\; t\in[0,\varepsilon),\ \;\varepsilon>0;\\


%


\dfrac{1}{(t^2+1)\arctg\,t},&\text{если }\; t\in[\varepsilon,\infty),


\end{cases}


\\


\|A\|=\int_0^{\infty} p(s)\,ds = \ln \frac{\pi}{2} -


\ln\arctg\varepsilon


\end{gather*}


может быть сколько угодно большой.


\end{note}





\begin{thebibliography}{9}





\bibitem{1}


Азбелев Н.В., Максимов В.П., Рахматуллина Л.Ф. {\it Введение в теорию


функционально-дифференциальных уравнений.} -- М.: Наука, 1991. -- 280 с.


\bibitem{2}


Исламов Г.Г. {\it Об оценке спектрального радиуса линейного положительного


вполне непрерывного оператора} // Функц.-дифференц. уравнения


и краев. задачи матем. физ. -- Пермь, 1978. -- С.\,119--122.


\bibitem{3}


Исламов Г.Г. {\it Об оценке сверху спектрального радиуса} //


Докл. РАН. -- 1992. -- Т.\,322. -- \No\,5. -- С.\,836--838.


\bibitem{4}


Азбелев Н.В., Рахматуллина Л.Ф. {\it Об оценке спектрального радиуса


линейного оператора в пространстве непрерывных функций} // Изв. вузов.


Математика. -- 1996. -- \No\,11. -- С.\,23--28.


\bibitem{5}


Кигурадзе И.Т., Шехтер Б.Л. {\it Сингулярные краевые задачи для


обыкновенного дифференциального уравнения второго порядка}


// Итоги науки и техн.


ВИНИТИ. Современ. пробл. матем. -- 1987. -- Т.\,30. --


С.\,105--201.


\bibitem{6}


Лабовский С.М. {\it О положительных решениях двухточечной краевой


задачи для линейного сингулярного


функционально-дифференциального уравнения} // Дифференц.


уравнения. -- 1988. -- Т.\,24. -- \No\,10. -- С.\,1695--1704.


\bibitem{7}


Красносельский М.А. , Вайникко Г.М., Забрейко П.П.


{\it Приближенное решение операторных уравнений.} --


М.: Наука, 1969. -- 455 с.





\end{thebibliography}





\vskip 2cm





\post{\it Пермский государственный}{\it Поступила}


\post{\it технический университет}{03.03.1998}





\label{end}


\end{document}


